\section{Pass manager mechanics}

\subsection{Proving invalidation instead of assuming it}

The new pass manager's preservation contract is easy to get subtly wrong. A
transform returns a \texttt{PreservedAnalyses} set, and it is tempting, or simply
habitual, to return \texttt{all()} and move on. If the transform actually changed
the IR, that leaves later passes reading stale cached analysis results. This is
exactly the class of mistake the project is meant to demonstrate command over, so
I did not want to trust my own eye on it.

The transforms here change instructions and values but never the control flow
graph, so the honest answer is to preserve the CFG set and nothing else. To make
sure that stays true, I wrote an invalidation test: it caches an opcode count,
runs the algebraic simplifier through a real function pass manager so the
invalidation machinery actually runs, then asks for the count again and asserts
it dropped. If the simplifier returned \texttt{all()}, the manager would hand back
the stale count and the test would fail.

I did not just assume the test was meaningful. I temporarily changed the
simplifier to return \texttt{PreservedAnalyses::all()}, rebuilt, and watched the
invalidation test fail; then I reverted and watched it pass. That experiment is
what makes the test trustworthy: it fails for the reason it claims to.

\subsection{The chain depth that could have looped forever}

The use-def chain depth metric walks operand dependencies looking for the longest
chain. In straight line code the dependency graph is acyclic, but a phi node in a
loop can make a value depend, transitively, on itself, which would send a naive
recursion into an infinite loop. I resolved it by treating phi nodes as chain
roots: the traversal stops at a phi and does not recurse into its operands. This
keeps the walk finite and gives the metric a clean meaning, the longest straight
line data dependency, with loop carried dependencies deliberately not counted. A
unit test on a loop kernel exists mostly to prove the traversal terminates at
all.
